\input{glyphtounicode}
\pdfgentounicode=1
\documentclass[11pt,a4paper]{article}
\usepackage{cmap}
\usepackage[utf8]{inputenc}
\usepackage[T1]{fontenc}
\usepackage[american]{babel}
\usepackage{lmodern}
\usepackage{microtype}
\usepackage[a4paper,top=2.5cm,bottom=2.5cm,left=2.5cm,right=2.5cm]{geometry}
\usepackage{amsmath,amssymb,amsthm}
\usepackage{booktabs}
\usepackage{tabularx}
\usepackage[colorlinks=true,linkcolor=blue,citecolor=blue,urlcolor=blue,unicode=true]{hyperref}
\usepackage{xurl}

\renewcommand*{\HyperDestNameFilter}[1]{en.#1}

\hypersetup{
  pdftitle={The Joshi ATS Series under the Bridge-Certificate Standard: A Targeted Substance Audit and an All-Place Rigidity Classification of Effective Normalization Weights},
  pdfauthor={Lukas Geiger},
  pdfsubject={Mathematics - Number Theory, Arithmetic Geometry},
  pdfkeywords={Arithmetic Teichmuller Space, abc conjecture, Bridge Certificate, Substance Audit, Rigidity Classification, Normalization Weights},
  pdfcreator={LaTeX with hyperref},
  pdfproducer={pdfTeX},
  bookmarksnumbered=true,
  pdfpagemode=UseOutlines,
  pdfstartview={FitH},
  unicode=true
}

\newtheorem{theorem}{Theorem}[section]
\newtheorem{lemma}[theorem]{Lemma}
\newtheorem{proposition}[theorem]{Proposition}
\newtheorem{corollary}[theorem]{Corollary}
\theoremstyle{definition}
\newtheorem{definition}[theorem]{Definition}
\theoremstyle{remark}
\newtheorem{remark}[theorem]{Remark}

\title{The Joshi ATS Series under the\\Bridge-Certificate Standard:\\[2pt]A Targeted Substance Audit and an\\All-Place Rigidity Classification of\\Effective Normalization Weights}
\author{Lukas Geiger\\\small Independent Researcher, Bernau, Germany\\\small ORCID: 0009-0005-7296-1534}
\date{October 2026 --- Version 0.3 (draft)}

\begin{document}
\maketitle

\begin{abstract}
This ongoing targeted audit separates coordinate bookkeeping, normalization
weights, and volume carriers in selected load-bearing steps of Joshi's ATS
series. The period-map argument requires consistent raw and normalized
logarithmic coordinates and a well-defined action on the period image.
Same-field all-place product-formula rigidity forces final effective
weights to be constant across places, but permits a global $j^2$ scaling;
its application across $L'/L$ and to the source's freeze remains conditional.
A local maximal-order hull inclusion is available under an actual module-hull
contract, while the ATS logarithm, Haar-degree, tensor and target-volume
bridges remain open. The printed ATS~III v4 corollary has absolute-value bars:
under a finite nonempty Tate-place sum with $0<|q_w|<1$ and positive
$\ell,\ell^*$, its left side is nonpositive and its right side positive.
This narrow sign inconsistency is separated from ATS~IV applications and
possible repairs. For global linear quantity families with defined reference
sums, rigidity gives at most a $y$-dependent scalar; classical height and
degree follow only with the stated normalization and domain identifications.
Measure and log-volume class membership and nonmembership are not established.
No conclusion about the truth of abc, the correctness of IUT, or the whole
ATS chain is claimed.
\end{abstract}

\clearpage
\tableofcontents
\clearpage

%=====================================================================
\section{Introduction}\label{sec:intro}

\subsection{Mandate and place in the series}

This series examines the published landscape around the abc conjecture
through a fixed auditing instrument. Part~1~\cite{part1} localized the
disputed step of Mochizuki's Corollary~3.12~\cite{mochizukiIUT,scholzestix};
Parts~2--3 developed the instrument: Part~2 introduced the \emph{bridge
contract} criterion~\cite{part2}, and Part~3 its certificate form~\cite{part3}
(the v1-F standard), which fix in advance what a transport argument between
two frameworks must deliver before its conclusion may be used;
Part~4~\cite{part4}
applied the certificate to Kirti Joshi's ATS series as a \emph{form audit}:
a pre-specified four-row test (\mbox{F-1}--\mbox{F-4}) of whether the chain states its transport law,
its compensation mechanism, its second mode, and its cross-container
comparability explicitly enough to be audited at all. The outcome of Part~4
was, deliberately, a statement about auditability, not about substance. A
tool supplement analyses the normalization identification
separately~\cite{gnc}, and a companion landscape paper situates the series
among the routes of the wider program~\cite{landscapeA}.

The present Part~5 is a \emph{targeted} substance audit of selected
load-bearing internal steps --- not a complete verification of the chain.
Its mandate, fixed before the work began, is the question left open by
Part~4: do the selected load-bearing steps of the published chain carry ---
that is, do the proofs written in the cited versions establish the
statements the chain uses at those sites? IUT imports, the existence base,
and steps outside the selected sites are outside the audit.

\subsection{Method}

The audit followed five phases recorded in a contemporaneous internal proof
register before the synthesis stage. Phase~0 mapped the
\emph{proof spine}: the minimal chain of results leading from the abc-type
inequality back to the existence of the deformation space, with an import
triage separating the chain's own results from statements quoted from the
IUT corpus. Phases~1--3 then verified, line-anchored against the exact
arXiv versions named in Section~\ref{sec:findings}, the sites the spine
prioritizes as heavily load-bearing: the full text of the normalization
mechanism, the base package (non-constancy theorem, renormalization
corollary, adelic assembly theorem), and the lower-bound core together with
its corollary translation. Phase~4 synthesized the findings into the typed
findings reported here; the reading question it left open is addressed by
the rigidity classification of Part~6, and Section~\ref{sec:case} reports
that classification under its named hypotheses.

Three disciplines from the earlier parts carry over unchanged. First,
\emph{verbatim anchoring}: every finding cites the exact version and, where
line-level extracts exist, the line range of the passage it concerns.
Second, \emph{typing}: findings are classified as formal bookkeeping
(verified or reconstructible), typed gaps (the text does not show what is
used --- with the missing step named), or documented reconstructions of our
own (marked as such, never silently supplied). Third, \emph{claim
discipline}: missing-proof findings are not promoted to falsity; the separate printed
sign test is stated only under its explicit hypotheses.

\subsection{Claim boundaries}\label{sec:claims-intro}

This paper reports the current state of an ongoing, targeted audit. The
same-field rigidity result is mathematical and conditional on its complete
place, field, and quantity-class hypotheses. Most source findings concern
missing implications, not counterexamples to a full theorem. Separately,
Section~\ref{sec:s7} identifies a narrow inconsistency in a printed sign
inequality under explicit finite Tate-place and normalization assumptions.
Neither kind of finding decides the truth of abc, IUT, or the whole ATS
chain. No new computation, machine-checked proof, or independent full-chain
validation is claimed. The audit's admissible claims may change when a
source revision, repair, or counterargument supplies the missing contracts.

The published v0.2 PDFs and the later, hash-bound bilingual source state are
distinct baselines. Intermediate work added Part~2 citation, table labels
and a response-route table, language-specific disclosure and revised layout.
The present scientific changes use the inventoried coordinate, field,
normalization, hull and sign corrections; the exact original v0.2 TeX was
not recovered. These provenance limits do not turn the intermediate source
into an independently verified copy of the published source.

\section{The Proof Spine of the Series}\label{sec:spine}

Phase~0 mapped the published chain backwards, from the abc-type conclusion to the
statement that deformation spaces of the required kind exist, using plain-text
extracts of the four versions listed at the beginning of
Section~\ref{sec:findings}. The map is complete \emph{within} those extracts and
truncated, not closed, wherever the chain calls on earlier papers of the
construction series that were not available for audit --- most consequentially at
the existence level, which rests on the first of them~\cite{joshiI}.

\subsection{The chain}

Table~\ref{tab:spine} lists the nodes lying immediately on the main chain, in the
order the proof text traverses them. The labels $S1,\dots,S22$ are our own
bookkeeping; the theorem numbers are the sources' own.

\begin{table}[htbp]
\centering
\footnotesize
\setlength{\tabcolsep}{4pt}
\caption{Load-bearing nodes, backwards from the abc-type conclusion. ``Proof''
counts extract lines between the \texttt{Proof.} marker and the end of the
argument: an indicator of checkability, not a measurement.}
\label{tab:spine}
\begin{tabularx}{\linewidth}{@{}l >{\raggedright\arraybackslash}X >{\raggedright\arraybackslash}X r@{}}
\toprule
 & Result & Rests on & Proof \\
\midrule
S1  & IV, Thm.~7.2.1 --- abc and arithmetic Szpiro over $\mathbb{Q}$ & S2 & none \\
S2  & IV, Thm.~7.1.1 --- Vojta inequality (declared main theorem) & S3, S9, S10 & 121 \\
S3  & IV, Thm.~6.1.1 --- First Main Bound & S4 alone & 33 \\
S4  & IV, Thm.~6.10.1 --- Second Main Bound (log-volume sandwich) & S5, S6, S7, S9 & $55{+}160$ \\
S5  & the middle equality of S4 (no theorem number) & ``same global normalization'' & 2 \\
S6  & IV, Prop.~6.10.9 / 6.10.10 / 6.10.12 --- upper bound & IUT~4, Thm.~1.10, steps (v)--(vii) & $2/1/2$ \\
S7  & III, Cor.~9.11.1.1 --- lower bound, translated & S8, by notational translation & none \\
S8  & III, Thm.~9.11.1 --- volume lower bound for $\widetilde{\Theta}$ & S15 and the hull and volume tools & 38 \\
S9  & IV, Thm.~5.7.1 --- existence of initial theta data & standard literature & 196 \\
S15 & III, Thm.~4.2.2.1 --- $\widetilde{\Sigma}_{L'}$ stable; $j^2$ scaling & ATS~II, ATS~II(1/2), Fargues--Fontaine & 15 \\
S16 & III, Cor.~4.2.2.4 --- forced renormalization & S15 and the ``global constraint'' & 2 \\
S17 & III, Thm.~4.6.1 --- global period map on $\widetilde{\Sigma}_{L'}$ & S18, S15, S16 & 3 \\
S18 & II(1/2), Thm.~5.10.1 --- product formula as period map & properties of $Y_L$ and arithmeticoids & 5 \\
S20 & II(1/2), Cor.~4.2.5 --- global Frobenius on $Y_L$ & S21 & 1 \\
S22 & II(1/2), Thm.~5.5.2 --- existence of arithmeticoids & Scholze; Kedlaya--Temkin; [Joshi 2021a] & 36 \\
\bottomrule
\end{tabularx}
\end{table}

In one sentence: S1 rests on the Vojta inequality S2, which rests on the First
Main Bound S3, which rests entirely on the log-volume sandwich S4; the
\emph{lower} half of that sandwich runs through S7 to S8, the \emph{upper} half is
taken over unchanged from IUT~4 (S6), and the \emph{middle equality} S5 is what
first joins the two halves into a single chain --- supported by the normalization
apparatus S15 to S18, which sits on the Frobenius construction S20/S21 and the
existence level S22. With the auxiliary nodes not shown, the map has 34 rows: 27
numbered nodes and 7 side rows, of which 15 lie immediately on the main chain.

\subsection{Import triage}

Each node was tagged by provenance: standard published literature; results the
series proves for itself; statements quoted from the IUT corpus; and purely
definitional passages. Keeping the third tag apart from the first is the
structurally important decision of the mapping, because the IUT corpus is not
background literature here but the artifact the chain sets out to replace. The
ledger of standard literature has 17 entries, 11 of them load-bearing in the sense
that a proof of the chain rests on them directly. The ledger of IUT imports has 13
entries, four of which are \emph{unexecuted full imports}, the entire proof
consisting of a pointer: ``This is the last equation on [Mochizuki, 2021d, Step
(v) on Page 658, Proof of Theorem~1.10] and its proof is all of step (v)''
(IV:3361--3362), with the parallel formulations for steps (vi) and (vii) at
IV:3372 and IV:3390--3391. These three propositions constitute the entire upper
half of the sandwich S4.

That fixes the boundary of this audit. Measured by load the IUT-quoted
propositions rank first; measured by novelty they rank last, being the source's
computation verbatim. A line-anchored audit there would be an audit of IUT~4, not
of the chain under examination, and lies outside the mandate of
Section~\ref{sec:claims-intro}. They are recorded and set aside, not verified.

\subsection{Proof text at the load-bearing nodes}

Of the 34 mapped rows, 21 proofs take one of three short forms: no proof block at
all (3), a proof consisting of an external citation (6), and a proof consisting of
one sentence of the form ``clear'', ``immediate'' or ``tautology'' (12). Twelve of
the 21 sit at nodes lying immediately on the main chain: ``The middle equality is
a tautology using the fact that the number fields provided by $y_0'$ and $y_0$ are
identically normalized'' (S5, IV:3272--3273); ``This is clear from
Theorem~4.2.2.1 and the global constraint placed by the requirement that the
product formula holds'' (S16, III:1613--1614); ``The proof is clear'' (S20,
II(1/2):1318).

The finding is descriptive. A two-line proof can be entirely sound, and several of
those listed are shown in Section~\ref{sec:findings} to carry exactly as written:
brevity here indexes \emph{checkability}, which an audit can address, not
correctness, which it cannot address in this form. The counts are counts of
extract rows, where page numbers stand as free-standing lines and displays break
across lines, so every length figure is an indicator and none carries an argument
below.

\subsection{What was selected for verification}

Phase~0 qualitatively prioritized the nodes by load borne, novelty, and
brevity of proof text, novelty being the discriminating factor for the
reason just given.

The ranking placed the middle equality S5 first: the only node joining the two
halves into one chain of inequalities, carrying no theorem number, dispatched in
the two lines quoted above, while the same paper states at three places that local
comparisons between $y_0$ and $\varphi(y_0)$ are precisely not available
(IV:3208--3211, IV:772--775, IV:816--819). Part~4 of this series treated S5 under
the transport-law row of the bridge certificate and resolved its status there ---
conditionally: the resolution holds on the bound and free readings of the
weight family laid open in Section~\ref{sec:fneu1} and not on the
place-selective reading, and Section~\ref{sec:case} reports how the
readings are classified under named hypotheses. S5 is therefore not
re-audited here.

The five remaining ranked nodes follow, preceded by the full-text question Part~4
had left open. \textbf{The normalization mechanism} --- what ``normalized
arithmeticoid'' operatively means --- comes first, because the reading it fixes is
presupposed by the three nodes after it. \textbf{S18} carries the entire
product-formula apparatus on a five-line proof whose justifying sentence states
that normalizations do \emph{not} carry over. \textbf{S16} asserts in two lines
that the $j^2$ scaling \emph{forces} a renormalization elsewhere, the precondition
for S17 and S18 being well defined at all. \textbf{S17} is marked by the author
himself as ``vital to establishing [Mochizuki, 2021c, Corollary~3.12]'' and as the
type of proposition ``claimed without proof'' in the works it competes with
(III:1838--1840), then proved in three lines. \textbf{S8} is the only
substantively independent lower bound in the chain, its decisive reduction
declared clear from the construction. \textbf{S7} is the node ATS~IV actually
cites: no proof block, presented as a notational translation of S8.

%=====================================================================
\section{Line-Anchored Findings}\label{sec:findings}

All findings below are anchored to four extracts; all four versions are
cited in their revisions of 24 February 2025. The line anchors refer to our
internal plain-text extracts and are finding aids, not a public
reproducibility standard; a provenance manifest (arXiv version, extraction
tool, checksums, line-numbering rule) is planned as a series supplement,
and until it exists the anchors should be read alongside the quoted
wording, which is verbatim.
\begin{itemize}
\item \textbf{IV} = ATS~IV, arXiv:2403.10430v2~\cite{joshiIV};
\item \textbf{III} = ATS~III, arXiv:2401.13508v4~\cite{joshiIII};
\item \textbf{II(1/2)} = ATS~II(1/2), arXiv:2305.10398v12~\cite{joshiIIhalf};
\item \textbf{II} = ATS~II (local prototype), arXiv:2303.01662v3~\cite{joshiII}.
\end{itemize}
An anchor of the form III:6840--6843 refers to lines of the plain-text extract,
not to pages of the published document. The extracts are lossy in one respect that
constrains what may be asserted: the II(1/2) extract loses Greek letters, so no
finding rests on where a $\lambda$ stands in a formula of that paper, whereas the
III and IV extracts preserve sub- and superscripts (control case III:1558) and the
$\xi$/$\Xi$ distinction. The present S7 formula is checked directly against the archived PDF,
whose absolute-value bars were lost by the extract; reconstructed formulas
are explicitly distinguished from that printed statement. One boundary recurs in Sections~\ref{sec:s8} and~\ref{sec:s7} and
is stated once here: Mochizuki's Corollary~3.12 is not part of the audited corpus,
so the claim that a statement of the chain \emph{is} that corollary
(III:6880--6881, IV:3197--3198) is not decidable here, and nothing below decides
it.

\subsection{The normalization mechanism}\label{sec:fneu1}

The normalization coordinate of ATS~II($\tfrac12$) \S\,5.3 is the positive
exponent converting the untilt absolute value to the fixed standard one.
Under the pointwise identity (5.3.3), direct substitution gives the classical
degree on $L^\times$ and height on $\mathbb P^n(L)$. This statement does not
automatically extend to $\mathbb P^n(\overline L)$ or line-bundle heights.
The compensation language in \S\,8 and ATS~III, Cor.~4.2.2.4, separately
requires an identified selection rule and field/place convention.

For a fixed number field $L$, all of its places, and real $x$-independent
effective weights satisfying the standard product formula for every
$x\in L^\times$, Part~6 gives $\beta_{y,v}=c_y$ across places. This leaves
a global scalar, possibly dependent on $y$; a constant global $j^2$
extension is admissible. What is excluded is unequal final weights under
these hypotheses, not every change of raw local normalization coordinates.
Tests on $L^\times$ at places of an extension $L'$ do not by themselves
establish the corresponding all-place theorem on $L'^\times$.

The source-bound question remains which rule and domain the operative
quantity uses. A pointwise-normalized height on $\mathbb P^n(L)$, a free
global scalar, a coordinate hyperplane, and a measure carrier are distinct
objects. No conclusion about the entire deformation apparatus follows
without the bridges in Sections~\ref{sec:s18}--\ref{sec:s8}.

\subsection{The period map: ATS~II(1/2), Theorem~5.10.1}\label{sec:s18}

The source asserts that $y\mapsto H_y$ is non-constant (ATS~II($\tfrac12$),
Thm.~5.10.1; extract II($\tfrac12$):2161--2243). Its local variation input is
available in ATS~II, Thm.~10.20.1(3), but the following coordinate and action
identifications must be separated from that input.

\paragraph{Coordinates and bookkeeping.}
Write raw logarithms as $u_v=\log|x_v|$ and normalized coordinates as
$z_v=\alpha_v u_v$. The source's degree is $\deg=\sum_v z_v$, whereas its
displayed weighted hyperplane is $H=\ker(\sum_v\alpha_v z_v)$.
Thus
\[
 \ker\Bigl(\sum_v z_v\Bigr)\ne
 \ker\Bigl(\sum_v\alpha_v z_v\Bigr)
 \quad\hbox{in general}.
\]
For $\alpha=(2,1)$ and $z=(1,-1)$ the first sum is $0$ and the second is $1$.
This is a small counterexample to an unqualified kernel identification.
It refutes neither the full theorem nor a consistently formulated
raw-coordinate reading.
A raw-log hyperplane $\ker(\sum_v\alpha_v u_v)$ is a separate consistent
contract; it must not be silently identified with the twice-weighted one.
The direct-sum support bookkeeping can be retained in its defined domain,
but the former blanket assertion that parts (1)--(5) carry is withdrawn.

\paragraph{Translations and the period image.}
A linear subspace $H$ is invariant under translation by $t$ exactly when
$t\in H$. Here this requires $\log_y(x)$ to lie in the actually chosen
hyperplane. The product formula for $\sum z_v$ alone does not establish
membership in $\ker(\sum\alpha_v z_v)$. An action on arithmeticoids descends
through a period map $F$ only if
\[
 F(y)=F(y')\ \Longrightarrow\ F(gy)=F(gy')
\]
for the relevant action elements $g$. That well-definedness is a separate
source obligation, not a consequence of local norm variation.

\paragraph{Non-proportionality and the remaining reconstruction.}
For a consistently defined weighted kernel, multiplying every coefficient
by the same nonzero scalar leaves the hyperplane unchanged. Non-constancy
of $\lambda_y$ therefore does not alone separate hyperplanes; a pair of
non-proportional coefficient families is needed. Reconstruction R-1 remains
conditional on H-REC: the varied non-archimedean local point can be inserted
while the other coordinates are held fixed. ATS~II, Thm.~10.20.1(3), supplies
the local variation input, not this recombination or the period-image action.
The archimedean reach remains unverified. No unconditional S18 closure is
claimed.

\subsection{The forced renormalization: ATS~III, Corollary~4.2.2.4}\label{sec:s16}

ATS~III, Cor.~4.2.2.4 (extract III:1600--1614), says that the $j^2$ scaling
at $V^{\mathrm{odd,ss}}$ forces a suitable renormalization at the remaining
places. This set concerns odd residue characteristic with bad semistable,
split-multiplicative reduction (ATS~III \S\,3.2(6)), not supersingular
reduction. Necessity, localization, existence, and uniqueness are separate
claims.

\paragraph{A baseline-relative necessity test (R-2).}
Suppose $S$ is a nonempty finite set of non-archimedean places of $L'$,
the tested global elements lie in $L^\times$, and the places restrict to
the specified usual absolute values on $L$. Relative to the additional
baseline of unchanged normalization coordinates away from $S$, the
scaled product has factor
\[
 \Bigl(\prod_{w\in S}|x|_w^{\mathrm{std}}\Bigr)^{j^2-1}.
\]
Choose a rational prime $r$ under one $w_0\in S$, so $r\in\mathbb Q^\times
\subset L^\times$. Then $\log|r|_w\le0$ on $S$ and is strictly negative at
$w_0$. For $j\ge2$ the displayed factor differs from $1$. This is a global
element in the correct test domain, not an arbitrary local uniformizer or
a principal-ideal generator potentially lying only in $L'$. Nonemptiness
of $S$ and the baseline are additional hypotheses; neither is inferred
merely from $\ell\ge5$.

\paragraph{Same-field rigidity and field descent (R-3).}
Let $b_v$ be the standard modulus family on one number field $K$, with
the canonical local degree factors included, so that
$\sum_v\log b_v(x)=0$ for every $x\in K^\times$. If
$\rho_{y,v}=b_v^{a_{y,v}}$ and $\lambda_{y,v}$ is the raw normalization
coordinate, the effective coefficient is $\beta_{y,v}=a_{y,v}\lambda_{y,v}$.
Rigidity constrains $\beta$, not $a$ and $\lambda$ separately.
For this modulus convention it gives $\beta=c\mathbf1$. If one instead
uses raw extended absolute values, the conclusion is proportionality to
their canonical degree-weight vector, not unweighted constancy. It requires
all places of $K$, including archimedean ones, real $x$-independent
coefficients and the full standard weighted product formula for every
$x\in K^\times$. Tests only on $L^\times$ across $w\in V_{L'}$ do not
imply constancy of each individual $L'$-weight. An explicit $L'/L$
descent/aggregation identity or the appropriate full $L'^\times$ formula
is needed. The arithmetic transfer is open; the same-field theorem is
not disputed.

\paragraph{What a genuine common freeze would imply.}
Under those same-field identifications and a nonempty freeze set, comparing
two solutions at the same $j$ with identical $a_{j,w}$ and frozen
$\lambda_{j,w}$ gives identical frozen $\beta_{j,w}$; constancy then gives
uniqueness if a solution exists. Nonemptiness is a stated prerequisite,
not a verified general assertion about the source. If, in addition, the
common reference family really satisfies
\[
 \rho_{j,w}=\rho_{1,w}^{j^2},\qquad
 a_{j,w}=j^2 a_{1,w},\qquad \lambda_{j,w}=\lambda_{1,w},
\]
at a frozen place, then algebraically
\[
 c_j=\beta_{j,w}=j^2\beta_{1,w}=j^2c_1.
\]
This admissible global extension does not violate rigidity. Its source
application still requires the common field, norm, place, and freeze
identifications; it yields no general verdict of projective emptiness.

\paragraph{The raw-coordinate freeze comparison.}
To avoid confusing two types of exponents, denote by $A_{j,v}$ the
coefficient vector in the fixed raw-coordinate comparison, not the local
absolute-value exponent $a_{j,v}$ above. In the inventoried comparison,
for $\varnothing\ne S\ne V$ and $j\ge2$, it has entries $\lambda_{1,v}$
on $S$ and $j^2\lambda_{1,v}$ outside $S$, and equals
$j^2\lambda_{\mathrm{can},j,v}$ under that comparison's canonical contract.
Consequently
\[
 H_{\mathrm{Freeze},j}=H_{\mathrm{can},j},
 \qquad H_{\mathrm{can},j}\ \hbox{need not equal}\ H_1.
\]
This conditional comparison does not certify the source's missing field
identification. Nor may a blanket projective-information conclusion be
drawn by identifying raw and effective vectors across different fields.
The literal finite-place convention of ATS~III and its possible
archimedean completion remain distinct open source contracts.

\subsection{The global period map: ATS~III, Theorem~4.6.1}\label{sec:s17}

The theorem has seven parts (III:1746--1817) and transports the apparatus of
Section~\ref{sec:s18} from a single arithmeticoid over $L$ to an $\ell^*$-tuple
over $L'$ constrained to Mochizuki's ansatz $\widetilde{\Sigma}_{L'}$. Its
complete proof: ``This is clear from [Joshi, 2023a, Theorem~5.10.1], and the
properties of $\widetilde{\Sigma}_{L'}$ established so far and detailed in
Theorem~4.2.2.1, Corollary~4.2.2.4'' (III:1819--1820).

Parts (1)--(3) and (5) cite the componentwise supplier, with separately
indexed target spaces (III:1749--1751). Their formal transport is conditional
on the consistent coordinate and action contracts of Section~\ref{sec:s18};
the former unqualified carry verdict is withdrawn. The tuple extension is
asserted in a source remark (II(1/2):2267), rather than proved there.
Part (4) explicitly invokes Corollary~4.2.2.4 (III:1791) and inherits the
field and normalization obligations of Section~\ref{sec:s16}. Part (7)
remains outside the passages examined.

\paragraph{The missing restriction step (T17-6).} Part~(6) asserts that
$z \mapsto (H_{y_1},\dots,H_{y_{\ell^*}})$ is non-constant \emph{on}
$\widetilde{\Sigma}_{L'}$, whereas the supplier, Theorem 5.10.1(7), asserts
non-constancy on all of $Y_L$ (II(1/2):2228--2233). But $\widetilde{\Sigma}_{L'}$
is a proper subset of the corresponding product, cut out by the sharp conditions
of Definition~4.2.2 (III:1460--1467), and non-constancy on a superset does not
imply non-constancy on a subset. None of the three sources cited in the proof
addresses that passage.

\paragraph{Reconstruction R-4: a conditional route.}
The source supplies stability of $\widetilde\Sigma_{L'}$ under global
Frobenius (Thm.~4.2.2.1(2), III:1536--1539) and the local norm relation
$|p|_{K_{y_{n-1}}}=|p|_{K_{y_n}}^{1/p}$ (ATS~II, Thm.~10.20.1(3)).
Neither statement alone proves non-constancy of the period map.
For the proposed route one must fix a common coordinate space and identify
the actual hyperplane coefficient vector $B_y$ in those coordinates.
One must then prove that a single admissible global move simultaneously
realizes the proposed local changes $\lambda_v\mapsto p_v\lambda_v$, that
these changes induce the claimed transformation of $B_y$ in the same
fixed coordinates, and that $B_{\varphi y}$ and $B_y$ are actually
non-proportional. Only for nonzero linear forms on that fixed space does
this imply $\ker B_{\varphi y}\ne\ker B_y$. Together with preservation of
$\widetilde\Sigma_{L'}$ and a well-defined period map, this would establish
non-constancy on that set. These are H-FROB obligations, not consequences
of varying residue characteristics or raw $\lambda$ coordinates alone.
The standard-point and archimedean caveats remain; the required global,
coordinate and coefficient identifications are not verified here.

\paragraph{A convention question that decides vacuity.} ATS~III writes the
logarithm and degree maps of parts (1) to (3) without the coefficients $\lambda_v$
(III:1760, III:1772, III:1781), and its standing convention declares all
arithmeticoids normalized. Read literally, the functional is the coordinate sum
and $H_{y_j}$ is the same hyperplane for every $z$, which would leave part (6)
without content rather than merely unproved; the statement becomes non-vacuous
only under the ATS~II(1/2) formulation with the $\lambda$-coefficients
(II(1/2):2212--2216), which ATS~III does not reproduce here. This is a text
finding, not an artifact of extraction.

\paragraph{Verdict.} \emph{Formal bookkeeping} for parts (1)--(3) and (5) remains conditional
on the corrected S18 coordinate and action contracts. \emph{Typed gap} at part (6): the proof as written
does not show non-constancy on $\widetilde{\Sigma}_{L'}$; the missing step is the
restriction from $Y_L$ to that subset. A second, inherited typed gap sits at part
(4); a third concerns the valuation convention, under whose literal reading part
(6) would have no content, the audited passages not fixing which reading governs.
\emph{Documented reconstruction}: R-4 is a conditional route: under H-FROB
(Section~\ref{sec:hyplabels}) it establishes non-constancy on
$\widetilde{\Sigma}_{L'}$, using sources the proof does not cite; without
an explicit global-to-local compatibility statement, the restriction step
remains open both in the source and in our reconstruction.

\subsection{The core of the lower bound: ATS~III, Theorem~9.11.1}\label{sec:s8}

ATS~III, Thm.~9.11.1 uses volume, whereas its template Thm.~9.9.1 uses a
supremum. Its box obligation must be assessed on the actual carrier, not
settled by observing that a single element has measure zero. In particular,
the claim that point membership cannot yield a positive-measure hull is
withdrawn.

\paragraph{A conditional algebraic hull inclusion.}
Suppose $A$ is a genuinely split finite local algebra, $\mathcal O_A^{\max}$
its specified maximal order, and $\mathcal O_{\mathrm{tensor}}$ the tensor
order with $\mathcal O_{\mathrm{tensor}}\subseteq\mathcal O_A^{\max}$.
Assume that the actual $\operatorname{hull}(S)$ contains $S$ and is stable
under multiplication by $\mathcal O_A^{\max}$, for example because it is
the corresponding maximal-order module hull. Then
\[
 x\in S\ \Longrightarrow\
 x\mathcal O_{\mathrm{tensor}}\subseteq
 x\mathcal O_A^{\max}\subseteq\operatorname{hull}(S).
\]
The last inclusion uses that module-hull hypothesis. A split algebra and
a maximal order alone do not establish it for a merely convex or
$\mathbb Z_p$-module hull. A positive full-dimensional box additionally
requires the componentwise nonvanishing and the actual local carrier.
No silent identification with the ATS target hull or $\widetilde\Theta$
is made.

\paragraph{Measure normalization and the outstanding volume contract.}
For a finite extension $E/\mathbb Q_p$, normalized additive Haar measure
$\mu_E(\mathcal O_E)=1$, and the absolute value
$|a|=|N_{E/\mathbb Q_p}(a)|_p^{c/[E:\mathbb Q_p]}$ with $c>0$ and
$|p|=p^{-c}$, change of variables gives
\[
 \mu_E(a\mathcal O_E)=|a|^{[E:\mathbb Q_p]/c}\quad(a\in E^\times).
\]
It is not generally $|a|$ under an arbitrary absolute-value convention.
The concrete logarithm image, ATS target hull, tensor and $\Gamma$ balance,
degree weights, collation and index layer must still be identified.
The abstract algebraic inclusion does not supply that volume theorem.

The finite-sum exponent identity used in the old comparison,
$\sum_{j=1}^{n}j^2/n^2\le n$ for $n\ge1$, remains an elementary
arithmetical fact. Its application to a measure lower bound requires the
specified Haar, degree, and carrier contracts; the former blanket
``exponent arithmetic above the gap is sound'' volume verdict is withdrawn.
The provenance and transport of the $\Xi$ objects through $\log_{BK}$
also remain separate source obligations (ATS~III \S\S\,9.8--9.11).

\subsection{The translated lower bound: ATS~III, Corollary~9.11.1.1}\label{sec:s7}

\paragraph{The printed statement.}
The archived ATS~III v4 PDF, printed p.~128, Cor.~9.11.1.1, actually has
absolute-value bars. With its notation, the displayed formula is
\[
 -\frac{1}{\ell^*}
 \left|\operatorname{LogVol}(\widetilde\Theta^{I}_{\mathrm{Mochizuki}})\right|
 \ \ge\ -\!
 \sum_{\substack{p,\ w\in V_p^{\mathrm{odd,ss}}\\
                  V_p^{\mathrm{odd,ss}}\ne\varnothing}}
 \log\left|q_w^{1/(2\ell)}\right|_{L'_w}.
\]
The summation is over the displayed prime blocks and their indicated
places; there are no additional real weights in this printed formula.
The old extraction lost the bars. The proposed atomic alternative
reading, the claim that the corollary needs bars supplied, and the former
exculpation through that alternative reading are explicitly withdrawn.

\paragraph{A narrow sign inconsistency, with hypotheses.}
Assume $\ell^*>0$, $\ell>0$, finite real log-volume, and a defined finite,
nonempty Tate-place sum with the usual Tate absolute-value normalization
$0<|q_w|_{L'_w}<1$. Each
$\log|q_w^{1/(2\ell)}|_{L'_w}$ is strictly negative. The left-hand side is
nonpositive and the right-hand side is strictly positive. Thus the printed
inequality cannot hold under these explicit assumptions. This is a local
sign inconsistency of that displayed assertion, not a verdict on the
whole ATS chain, IUT, or abc. Empty sums, divergent expressions, or
different norm conventions are not silently covered.

\paragraph{What this does not establish.}
A lower bound $\mathrm{Vol}\ge R$ with $0<R<1$ does not imply
$\mathrm{Vol}\le1$ or $\operatorname{LogVol}\le0$; $\mathrm{Vol}=2$ is a
counterexample to that implication. The source's separate local
nonpositivity assertion must not be substituted for a proved volume
contract. The criticism of the printed ATS~III form, its use in ATS~IV,
and any corrected lower-bound statement require separate comparisons.
A repaired sign convention alone would not close the hull/measure
obligations of Section~\ref{sec:s8}. H-SIGN below records the actual
printed bars and the sign-test assumptions, not an alternative reading.

\section{Typed Gap Catalogue and Verdict Table}\label{sec:catalogue}

\subsection{Verdict table}

\begin{table}[htbp]\centering\small
\caption{Verdict by operation; open bridges are not whole-chain verdicts.}
\label{tab:verdict}
\begin{tabularx}{\textwidth}{@{}>{\raggedright\arraybackslash}p{0.32\textwidth} >{\raggedright\arraybackslash}X@{}}
\toprule Operation & Current scope \\\midrule
S18 coordinates/action & direct-sum bookkeeping only in its domain; weighted kernel and period-image action separate \\
S18 non-proportionality & R-1 conditional on consistent coordinates and H-REC; archimedean reach open \\
S16 normalization & same-field rigidity valid; $L'/L$, nonempty freeze and norm identifications open; global $j^2$ allowed \\
S8 hull/volume & local inclusion conditional on actual module-hull stability; ATS Haar/tensor/target-volume contract open \\
S7 printed corollary & absolute-value bars present; narrow sign inconsistency under H-SIGN; IV and repairs separate \\
S17 restriction & source restriction step open; R-4 conditional on H-FROB \\
\bottomrule\end{tabularx}\end{table}

\subsection{Documented reconstructions}

R-1 is the non-proportionality route under consistent coordinates and H-REC;
R-2 is the baseline-relative necessity test with a global rational prime;
R-3 is same-field rigidity, whose source application requires field descent;
R-4 is the conditional global-period restriction route under H-FROB.
The local hull inclusion is conditional on its actual maximal-order module
contract. None closes the complete ATS volume or period chain.

\subsection{A comparative observation}\label{sec:symmetry}

Within the comparison corpus of Part~1~\cite{part1} and the present audit,
both programs contain an unexpanded containment-type transition at a
structurally load-bearing site: the program the chain seeks to replace
asserts its decisive step with ``follows formally''
(cf.~Part~1), the chain under audit asserts its box containment with
``clear from the construction'' (Section~\ref{sec:s8}). We do not assert
identity of the objects, equivalence of the missing implications, or a new
audit of IUT here; the observation is a comparison of \emph{assertion form}
at the respective critical sites, no more. The difference in form otherwise
remains as Part~4 recorded it --- the chain under audit states a theorem
with a written-out proof and nameable repair obligations. The elementary
exponent identity is distinct from its still-open measure application. Neither program is refuted by this observation.

%=====================================================================
\section{The Rigidity Classification of the Normalization Weights}\label{sec:case}

This section applies the rigidity results of Part~6~\cite{part6} to the
normalization center isolated in Section~\ref{sec:fneu1}. The statements
below are implications under explicitly named hypotheses, not pairwise
disjoint readings of the source; the hypotheses are defined once here and
used consistently.

\subsection{Hypothesis labels}\label{sec:hyplabels}

\begin{description}
\item[H-PLACES.] One fixed number field $K$ and its full set of places,
including archimedean ones, with an identified standard absolute-value family.
An abbreviation is admissible only when the omitted factors are demonstrably
retained. Tests on $L^\times$ at $L'$-places require a separate descent bridge.
\item[H-PF.] Real, $x$-independent final effective weights satisfy the
full standard product formula for every $x\in K^\times$, without suppressed
place, fiber, or correction terms.
\item[H-5.3.3.] The pointwise identity (5.3.3) applies at every specified
place and identifies the effective coefficients with $\beta_{y,v}=1$.
Its use on $L^\times$ and $\mathbb P^n(L)$ does not establish extension-field
or line-bundle transports.
\item[H-EFF.] The quantity has a global \emph{linear} representation
\[
 Q_y=\sum_v\beta_{y,v}q_v,\qquad Q_{\mathrm{ref}}=\sum_vq_v,
\]
where $q_v$ are $y$-independent and both sums are defined with compatible
index/order conventions and sufficient convergence for the stated scalar
factorization. Finite sums or absolutely convergent real sums suffice.
No arbitrary nonlinear factorization is included. Identifying
$Q_{\mathrm{ref}}$ with $Q_{\mathrm{cl}}$ is an additional classical bridge.
\item[H-ID.] The operative height/degree is the specifically defined
quantity on $\mathbb P^n(L)$ or $L^\times$, with no uncontrolled additional
$y$-data. Algebraic extensions and line bundles require separate bridges.
\item[H-REC.] The required varied local point can be recombined into an
admissible arithmeticoid while the other specified coordinates stay fixed;
consistent log/degree/hyperplane coordinates and any period-image action
are separate obligations.
\item[H-FROB.] The global Frobenius move of R-4 realizes the required
local transformations simultaneously and stays in $\widetilde\Sigma_{L'}$.
It identifies the actual nonzero hyperplane forms in a common fixed
coordinate space, proves their stated transformation and actual
non-proportionality there, and supplies the well-defined period-map
contract. Raw normalization-coordinate variation alone does not suffice.
\item[H-SIGN.] The literal printed absolute-value formula of
Cor.~9.11.1.1, with $\ell,\ell^*>0$, finite real log-volume and a finite,
nonempty usual Tate-place sum with $0<|q_w|_{L'_w}<1$, is subjected to the
sign test in Section~\ref{sec:s7}. This is not an alternative convention.
\end{description}
H-5.3.3 implies H-PF with $\beta=1$ in the defined same-field standard
contract. H-PF alone yields only $\beta=c_y\mathbf1$, not $y$-independence.
H-EFF concerns a defined linear quantity family; it is not a field or place
identification and does not classify measures automatically.

\subsection{The readings of the weight family}

The normalization analysis of Section~\ref{sec:fneu1} identified three
readings of the weight family $(\lambda_{y,v})$ in the audited text, and
found that the text does not select among them: a \emph{bound} reading, on
which the normalization coordinate is determined by the Artin--Whaples
standard normalizations~\cite{artinwhaples} (``\emph{the} normalization
coordinate'' of \S\,5.3); a \emph{free} reading, on which the weights are any
real family satisfying the weighted product formula; and a
\emph{place-selective} reading, on which a rule assigns genuinely
place-selective weights while the product formula continues to hold ---
this last being what the chain's \emph{use} of the weights requires. These
readings are not disjoint classes of solutions: the bound tuple is itself a
product-formula solution, and a place-selective solution, if one existed,
would be one too. They are alternatives for what the source \emph{means};
which of them the source intends is a question of text interpretation that
the mathematics below does not decide. What the mathematics decides is
which of these readings are \emph{realizable} by exponent families under
the named hypotheses.

\subsection{Applicability gate: range of places}\label{sec:gate}

The rigidity theorem of Part~6 applies to a full all-place product formula.
Before it can be applied to the audited passages, the place-range conflict
recorded in Section~\ref{sec:s16} must be confronted: ATS~II($\tfrac12$)
introduces $V_L$ as all places, archimedean ones included, whereas ATS~III
\S\,3.2(5), read literally, indexes only non-archimedean valuations and runs
its product formula over that range. The audited text does not reconcile
these conventions, and power-class equivalence of the local absolute values
does not substitute for the missing places. Therefore: (i) if ATS~III
suppresses but retains the archimedean factors, and the full same-field
test domain and other H-PF hypotheses are supplied, the rigidity theorem
applies to that completed effective-weight vector; (ii) if
ATS~III's product formula is literally a non-archimedean-only constraint,
the theorem does not apply as stated, and the status of that passage
remains open pending a separate theorem or an explicit archimedean
completion. In the special case where the finite-place compensation is
combined with fixed standard archimedean weights, the combined all-place
vector falls under H-PLACES and H-PF, and rigidity then forces the finite
weights to a single global value; but this combined formulation must be
stated by the source, not inferred from notation. Everything in the
remainder of this section is asserted under H-PLACES.

\subsection{The classification from Part~6}

Part~6 v0.3~\cite{part6} supplies same-field rigidity and unchanged
finite-fiber redistribution. Under H-PLACES and H-PF, $\beta_{y,v}=c_y$.
Under H-EFF, with its actual convergence contract,
\[
 Q_y=c_y Q_{\mathrm{ref}}.
\]
Only with the additional identification $Q_{\mathrm{ref}}=Q_{\mathrm{cl}}$
may this be written as $c_yQ_{\mathrm{cl}}$. Under H-5.3.3 the coefficients
are $1$; direct substitution yields the classical degree on $L^\times$
and height on $\mathbb P^n(L)$. Neither algebraic-closure points nor
ATS~III line-bundle heights inherit this merely by notation.

\begin{table}[htbp]\centering\small
\caption{Rigidity consequences with explicit domains and quantity contracts.}
\label{tab:rigidity-consequences}
\begin{tabularx}{\textwidth}{@{}>{\raggedright\arraybackslash}p{0.32\textwidth} >{\raggedright\arraybackslash}X@{}}
\toprule Contract & Admissible conclusion \\\midrule
H-PLACES, H-PF & constant final weights; global $j^2$ remains admissible \\
H-EFF as well & $Q_y=c_yQ_{\mathrm{ref}}$; classical identification separate \\
H-5.3.3 and H-ID & classical degree on $L^\times$, height on $\mathbb P^n(L)$ \\
Measures, hulls, log-volume & membership and nonmembership in H-EFF unclassified; no rigidity conclusion \\
\bottomrule\end{tabularx}\end{table}

For a stabilized expression on a specified nonempty $L^\times$ family,
H-5.3.3 gives the same classical value for every defined constituent only
under the stated source normalization; its supremum then has that value.
Different domains or added carrier data require separate bridges. A general
supremum is not silently treated as a linear H-EFF quantity, nor is a
scalar pulled through it without the necessary sign/domain assumptions.
The coordinate variation of $H_y$ alone is no proof of metric variation.

\subsection{Quantity-class bridge: what follows and what does not}\label{sec:bridge}

The quantity-class bridge is a domain-sensitive obligation, not a
generic transfer of coordinate variation. A degree on $L^\times$ or
height on $\mathbb P^n(L)$ is classified only by its actual H-EFF/H-ID
substitution. A stabilized expression requires its own domain and
normalization contract. A hyperplane family alone gives no metric
variation. Log-volume, hull, Haar measure and box inclusion have neither
membership nor nonmembership in H-EFF established here.
\begin{table}[htbp]\centering\small
\caption{Quantity-class obligations from weights to operative carriers.}
\label{tab:quantity-class-bridge}
\begin{tabularx}{\textwidth}{@{}>{\raggedright\arraybackslash}p{0.32\textwidth} >{\raggedright\arraybackslash}X@{}}
\toprule Carrier & Required bridge \\\midrule
$L^\times$ degree / $\mathbb P^n(L)$ height & defined linear substitution and classical reference identification \\
Stabilized expression & specified nonempty domain, normalization and any supremum/scalar conditions \\
$H_y$ & consistent coordinates and well-defined period action; no metric conclusion by itself \\
Hull, measure, log-volume & actual carrier, Haar-degree, tensor/\,$\Gamma$ and volume contract; class status unclassified \\
\bottomrule\end{tabularx}\end{table}
Thus Part~6 excludes unequal final place weights under H-PLACES/H-PF;
it does not decide the entire height, measure or log-volume chain without
the identified quantity/domain bridges. No downstream-node failure is
inferred from the classification alone.

\section{Response-Verification Template and Conclusion}\label{sec:template}

\subsection{The decision tree}

The findings of this paper concern the published text state at the audit
dates given. A future revision of the ATS series, or a reply by its author,
is the expected and welcome next event; the following decision tree, fixed
here in advance, records the \emph{alternative} routes by which such a
document could change the findings. The routes are not cumulative
requirements: each branch applies only if that route is chosen, and a
legitimate refutation of this audit itself is a further admissible outcome
on every branch.

Table~\ref{tab:response-routes} states the route logic before the prose
checklist.

\begin{table}[htbp]
\centering
\small
\caption{Response routes and what each would need to establish.}
\label{tab:response-routes}
\begin{tabularx}{\textwidth}{@{}>{\raggedright\arraybackslash}p{0.25\textwidth} >{\raggedright\arraybackslash}p{0.28\textwidth} >{\raggedright\arraybackslash}X@{}}
\toprule
Route & What would change & Minimum showing \\
\midrule
Misidentification & H-ID no longer attaches the classification to the
load-bearing quantity & identify the operative quantity and its additional
$y$-data \\
Exponent-class repair & the source remains within H-EFF but repairs the
missing all-place data & settle H-PLACES, give a selection rule, and track
any global scalar $c_y$ \\
Outside-H-EFF carrier & the rigidity theorem no longer obstructs the chosen
carrier & define the measure or log-volume carrier and prove that it carries
the needed variation \\
Lower-bound route & normalization is no longer the decisive step & prove
full box containment or supply an alternative lower-bound theorem \\
Text/citation completion & documentary closures become reconstructible &
write the R-1/R-4 closures with sources \\
\bottomrule
\end{tabularx}
\end{table}

\begin{enumerate}
\item \textbf{Misidentification route.} Show that the operative
  load-bearing height is not the height of Definition~8.2.1, or that it
  contains additional $y$-data --- i.e.,\ refute H-ID for the audited
  passages. The classification of Section~\ref{sec:case} then no longer
  attaches to the load-bearing quantity.
\item \textbf{Repair within the exponent class.} If this route is chosen:
  settle the place range of Section~\ref{sec:gate} (H-PLACES), state an
  explicit selection rule or a proof of choice-independence for the forced
  renormalization (the well-definedness obligation of
  Section~\ref{sec:s16}), verify the rule against the one-dimensional
  solution space at $L = \mathbb{Q}$, and track the effect of a possible
  global scalar $c_y$ through every used expression.
\item \textbf{Carrier outside the effective-weight class.} If the
  $y$-dependence is to be carried by a quantity outside H-EFF --- for
  instance at the level of measures or log-volumes --- define that
  quantity, give its address in the chain, and prove that it carries the
  required variation; the rigidity theorem does not obstruct this route,
  and nothing in this paper closes it.
\item \textbf{Lower-bound route.} Independently of normalization: either
  prove the full box containment used in Thm.~9.11.1
  (Section~\ref{sec:s8}), or supply an alternative lower-bound theorem with
  a complete derivation.
\item \textbf{Text and citation completion.} After any of the above:
  write the reconstructible closures of Sections~\ref{sec:s18}
  and~\ref{sec:s17} into the text with their sources. This is a
  documentation obligation, not a mathematical precondition of the other
  routes.
\end{enumerate}

\subsection{Conclusion}

The current targeted audit retains same-field product-formula rigidity
and the distinction between non-constancy and non-proportionality. It
separates raw and normalized coordinates, the period-image action,
field/place descent and actual freeze identifications. Global $j^2$
scalings are admissible; unequal final weights are excluded only under
the complete same-field all-place hypotheses. A conditional algebraic
module-hull inclusion is available, while the ATS logarithm, Haar-degree,
tensor and target-volume contracts remain open. The actual printed
absolute-value corollary has the narrow sign inconsistency stated in
Section~\ref{sec:s7}; its use in IV and any repaired statement are separate.
Linear quantity-class factorization requires defined reference sums and
convergence, and classical heights require their stated domains.
Measure membership and nonmembership remain unclassified.

This is the ongoing research project's present scope, not an independent
verification of the whole ATS chain or a proof/refutation of abc or IUT.
The response routes above remain open to a source revision, explicit
bridge, or a counterargument to this audit itself. No new computation or
machine-checked proof is supplied.

%=====================================================================
\section*{AI Disclosure}
AI systems assisted source work, mathematical discussion, drafting,
translation and typesetting. Historical adversarial reviews and their
limited scopes are recorded in the research files. The present corrective
preparation uses the inventoried findings and a narrow original-page check;
it introduces no new numerical experiment or machine-checked proof.
Independent acceptance of this candidate is not claimed by its preparation.


%=====================================================================
\begin{thebibliography}{99}
\raggedright

\bibitem{artinwhaples}
E.~Artin and G.~Whaples,
\emph{Axiomatic characterization of fields by the product formula for
valuations},
Bull.\ Amer.\ Math.\ Soc.\ \textbf{51} (1945), no.~7, 469--492.
DOI: 10.1090/S0002-9904-1945-08383-9

\bibitem{joshiI}
K.~Joshi,
\emph{Construction of Arithmetic Teichmuller Spaces I},
arXiv:2106.11452v4 [math.AG], 2021 (revised 24 February 2025).

\bibitem{joshiII}
K.~Joshi,
\emph{Construction of Arithmetic Teichmuller spaces II: Proof of a local
prototype of Mochizuki's Corollary~3.12},
arXiv:2303.01662v3 [math.AG], 2023 (revised 24 February 2025).

\bibitem{joshiIIhalf}
K.~Joshi,
\emph{Construction of Arithmetic Teichmuller Spaces II$\frac{1}{2}$:
Deformations of Number Fields},
arXiv:2305.10398v12 [math.AG], 2023 (revised 24 February 2025).

\bibitem{joshiIII}
K.~Joshi,
\emph{Construction of Arithmetic Teichmuller Spaces III: A `Rosetta Stone'
and a proof of Mochizuki's Corollary~3.12},
arXiv:2401.13508v4 [math.AG], 2024 (revised 24 February 2025).

\bibitem{joshiIV}
K.~Joshi,
\emph{Construction of Arithmetic Teichmuller Spaces IV: Proof of the
abc-conjecture},
arXiv:2403.10430v2 [math.AG], 2024 (revised 24 February 2025).

\bibitem{mochizukiIUT}
S.~Mochizuki,
\emph{Inter-universal Teichm\"{u}ller Theory I: Construction of Hodge
Theaters}, Publ.\ Res.\ Inst.\ Math.\ Sci.\ \textbf{57} (2021), no.~1/2,
3--207;
\emph{II: Hodge--Arakelov-Theoretic Evaluation}, \emph{ibid.}, 209--401;
\emph{III: Canonical Splittings of the Log-Theta-Lattice}, \emph{ibid.}, 403--626;
\emph{IV: Log-Volume Computations and Set-Theoretic Foundations},
\emph{ibid.}, 627--723.

\bibitem{scholzestix}
P.~Scholze and J.~Stix,
\emph{Why abc is still a conjecture},
unpublished manuscript, version of 16 July 2018.
Available at
\url{https://www.math.uni-bonn.de/people/scholze/WhyABCisStillaConjecture.pdf}

\bibitem{part1}
L.~Geiger,
\emph{The Gap in IUTchIII, Corollary~3.12: From Attempted Repair to
Structural Diagnosis},
Zenodo, 2026.
Concept DOI: 10.5281/zenodo.19960781
(latest version v6.3, 12 August 2026, DOI: 10.5281/zenodo.21899747).

\bibitem{part2}
L.~Geiger,
\emph{The Bridge-Contract Criterion: A Form Audit for Transport Arguments,
with an Instantiation in the abc Literature},
Zenodo, 2026.
Concept DOI: 10.5281/zenodo.21960476
(version 0.2, 16 August 2026, DOI: 10.5281/zenodo.21960477).

\bibitem{part3}
L.~Geiger,
\emph{The Bridge Certificate: A Line-Level Audit Instrument for
Inter-Theoretic Transfer Arguments, with Four Applications to the abc
Literature (DRAFT)},
Zenodo, 2026.
Concept DOI: 10.5281/zenodo.21925803
(version 0.2, 16 August 2026, DOI: 10.5281/zenodo.21960958).

\bibitem{part4}
L.~Geiger,
\emph{The Bridge Certificate Applied to the Joshi ATS Series: A Ledger
Audit},
Zenodo, 2026.
Concept DOI: 10.5281/zenodo.21951155
(version 0.2, 15 August 2026, DOI: 10.5281/zenodo.21951156).

\bibitem{part6}
L.~Geiger,
\emph{What the Product Formula Allows: A Rigidity Theorem and an
Effective-Weight Test for the ATS Height Subchain},
Zenodo, 2026.
Concept DOI: 10.5281/zenodo.21952486
(version 0.3, 2 October 2026, DOI: 10.5281/zenodo.23093243).

\bibitem{gnc}
L.~Geiger,
\emph{Global Normalization in Joshi's Arithmetic Teichm\"{u}ller Theory: What
the Identification Carries --- and What Carries the Theorem},
Zenodo, 2026.
Concept DOI: 10.5281/zenodo.21924253
(version 1.1, 13 August 2026, DOI: 10.5281/zenodo.21924254).

\bibitem{landscapeA}
L.~Geiger,
\emph{The abc Landscape: Obstructions, Failed Routes, and HCT Diagnostics
(DRAFT)},
Zenodo, 2026.
Concept DOI: 10.5281/zenodo.21916900
(latest version 1.1, 13 August 2026, DOI: 10.5281/zenodo.21924338).

\end{thebibliography}

\end{document}
